\chapter{Session layer}\label{sec:session}

The Cyphal session layer provides the named publish-subscribe abstractions that
applications use directly.
It sits above the transport layer (chapter~\ref{sec:transport}), treating it as a narrow,
unreliable, deduplicated, unordered message courier that carries multicast and unicast
transfers between nodes.
Every feature that requires more than that contract -- naming, subject-ID allocation,
reliable delivery, remote procedure calls, streaming, service discovery -- is
constructed by the session layer from a small number of session-layer headers and
state machines described in this chapter.

\section{Function and position in the stack}\label{sec:session_function}

The session layer provides, on top of the transport contract defined in
section~\ref{sec:transport} and summarized below:

\begin{itemize}
    \item A \emph{named topic} abstraction with decentralized subject-ID allocation,
    replacing the static subject-ID tables used by Cyphal v1.0.
    \item Reliable one-way message delivery with acknowledgements and retransmission.
    \item Remote procedure calls with optional response streaming, built as a pattern on
    top of ordinary pub-sub transfers rather than as a separate service-ID construct.
    \item Pattern-based subscriptions that double as a service discovery mechanism.
    \item A gossip-based consensus protocol that keeps every node's view of the
    topic-to-subject-ID map consistent with every other node's view.
\end{itemize}

The transport layer is required only to provide unreliable, deduplicated, unordered
delivery of integrity-checked messages; segmentation and reassembly of large payloads;
and participant discovery as a side effect of joining the broadcast subject.
These requirements are stated in full in section~\ref{sec:transport};
the session layer does not depend on anything else.

\begin{remark}
    A consequence of this division is that the session layer is portable across
    transports without modification.
    A new Cyphal transport is made interoperable with the rest of the ecosystem by
    providing the transport contract only; every session-layer behavior described in
    this chapter is expected to work on top of it without change.
\end{remark}

\section{Model and core objects}\label{sec:session_model}

The session layer operates on the following objects.
Each object is described in terms of the state it carries and the wire messages it
exchanges; no application programming interface is implied.

\begin{description}
    \item[Node] A participant in a Cyphal network.
    A node owns a nonempty \emph{home name} that should be unique in the network,
    a \emph{namespace} prefix used when resolving topic names,
    and an optional \emph{remap table} that rewrites topic names at resolution time.
    A node exposes publishers and subscribers and participates in the allocation
    consensus protocol by exchanging gossip messages with other nodes.

    \item[Topic] A named publish-subscribe channel.
    A topic is identified on the wire by a 64-bit \emph{topic hash} and a
    transport-visible \emph{subject-ID}; the mapping between them is maintained by
    the allocation protocol described in section~\ref{sec:session_crdt}.
    Locally, a topic carries a small CRDT state record $(\mathit{hash},\
    \mathit{evictions},\ \mathit{age})$ whose semantics are defined in
    section~\ref{sec:session_crdt}.

    \item[Publisher and subscriber] Per-topic endpoints through which the application
    emits and consumes messages. They are named here only because their wire
    behavior differs (publishers emit messages, subscribers emit acknowledgements);
    they are not otherwise observable to other nodes.
    A single topic may have several local publishers and subscribers simultaneously.

    \item[Future] A conceptual handle representing the progress of an asynchronous
    operation -- a reliable publication waiting for an acknowledgement, a request
    waiting for responses, or a subscriber waiting for an arrival.
    Futures are used in this chapter only to describe observable state transitions
    of the protocol; the wire format does not reference them.
\end{description}

\section{Topic names}\label{sec:session_names}

A \emph{topic name} is a string of printable ASCII characters used by the application
to identify a publish-subscribe channel.
The session layer carries names only in gossip and scout messages; every other
message identifies topics by their 64-bit hash (section~\ref{sec:session_identity}).

\subsection{Character set and normalization}

A valid name consists of characters in the ASCII range $[33, 126]$; the space
character (ASCII $32$) and all control characters are excluded.
Names are case-sensitive.
The forward-slash character \texttt{/} (ASCII $47$) acts as the \emph{separator}
between name segments.

A \emph{normalized} name contains no leading separator, no trailing separator, and
no consecutive separators. Names are normalized before hashing and before being sent
on the wire. The maximum length of a normalized name is 200 bytes.

\subsection{Pattern tokens}

Subscribers may register under pattern names that contain the following special
tokens; publishers may not. A publisher whose resolved name contains any pattern
token shall be rejected by the session layer.

\begin{description}
    \item[\texttt{*} (ASCII $42$)] Matches exactly one name segment.
    For example, the pattern \texttt{sensors/*/data} matches \texttt{sensors/imu/data}
    and \texttt{sensors/gps/data} but not \texttt{sensors/imu/raw/data}.

    \item[\texttt{>} (ASCII $62$)] Matches zero or more trailing segments.
    Must appear as the last token in the pattern.
    For example, the pattern \texttt{sensors/>} matches \texttt{sensors},
    \texttt{sensors/imu}, and \texttt{sensors/imu/raw/data}.
\end{description}

Pattern tokens are not permitted in intermediate positions outside of these forms;
for example, \texttt{sen*ors} is a literal name, not a pattern.

\subsection{Pin suffix}

A name may be pinned to a specific subject-ID using the suffix
\texttt{\#\textit{NNN}}, where \texttt{\#} is ASCII $35$ and \textit{NNN} is a
decimal number in the range $[0, 8191]$ without leading zeros.
Pinning is permitted only on verbatim names; a pattern name shall not carry a
pin suffix.
The pin suffix is stripped from the name before hashing and normalization, so a
pinned topic and its verbatim counterpart share the same topic hash but request a
different subject-ID (section~\ref{sec:session_identity}).

\subsection{Home and namespace expansion}

A name may begin with the \emph{home} token \texttt{\textasciitilde} (tilde,
ASCII $126$), which is replaced at resolution time with the node's home name.
A name beginning with \texttt{/} is \emph{absolute} and is used as-is, ignoring
the namespace prefix. All other names are relative and are prefixed with the
node's namespace at resolution time.

\subsection{Remapping}

A node may carry a \emph{remap table} of \texttt{\textit{from}}$\to$\texttt{\textit{to}}
pairs that rewrite names before namespace and home expansion.
A matching rule replaces the input name with its \texttt{\textit{to}} string and then
feeds the result into the rest of the resolution pipeline.
If the matched rule's \texttt{\textit{to}} string carries a pin suffix, that pin is
honoured; any pin supplied in the user's original name is discarded.
Software implementations should accept the remap table from the environment variables
\texttt{CYPHAL\_NAMESPACE} and \texttt{CYPHAL\_REMAP} to allow integrators to adjust
names without modifying the application.

\subsection{Resolution pipeline}

The following pipeline is normative. Given a user-supplied name, a node resolves it
as follows, in order:

\begin{enumerate}
    \item \textbf{Remap} the input name through the remap table, if any.
    \item \textbf{Expand} the home token \texttt{\textasciitilde} if present, or
    prefix the namespace if the name is relative and the absolute form is not used.
    \item \textbf{Extract} any pin suffix, recording the pinned subject-ID separately.
    \item \textbf{Normalize} the remaining name by collapsing consecutive separators
    and trimming leading and trailing separators.
\end{enumerate}

The result is the normalized resolved name (used for hashing and for inclusion in
gossip and scout messages) together with the optional pinned subject-ID.
Resolution fails if the normalized name is empty, exceeds 200 bytes, or contains a
character outside the permitted set.

\begin{remark}
    The intended usage pattern is that applications hardcode verbatim names and let
    integrators introduce patterns and pinning through the remap table at deployment
    time. This keeps the application code free of site-specific routing decisions.
\end{remark}

\subsection{Resolution examples}

The following examples illustrate the resolution pipeline when no remap table
is installed. \texttt{NAMESPACE} and \texttt{HOME} are per-node configuration;
\texttt{RESOLVED} is the normalized name that appears on the wire;
\texttt{PIN} is the pinned subject-ID, if any;
\texttt{VERBATIM?} is \emph{yes} if the resolved name has no pattern tokens.

\begin{CyphalSimpleTable}[wide]{Name resolution without remap}{|l l l l r l|}
    Name                   & Namespace   & Home    & Resolved           & Pin    & Verbatim? \\
    \texttt{foo/bar}       & \texttt{ns1}   & \texttt{me} & \texttt{ns1/foo/bar}       & --       & yes \\
    \texttt{\textasciitilde/foo/bar} & \texttt{ns1} & \texttt{me} & \texttt{me/foo/bar} & -- & yes \\
    \texttt{/foo//bar/}    & \texttt{ns1}   & \texttt{me} & \texttt{foo/bar}           & --       & yes \\
    \texttt{foo/bar/}      & \texttt{\textasciitilde/ns1} & \texttt{me} & \texttt{me/ns1/foo/bar} & -- & yes \\
    \texttt{foo\#123}      & \texttt{ns1\#456} & \texttt{me} & \texttt{ns1/foo}         & $123$    & yes \\
    \texttt{foo/\#123}     & \texttt{ns1\#456} & \texttt{me} & \texttt{ns1/foo}         & $123$    & yes \\
    \texttt{*/foo}         & \texttt{ns1}   & \texttt{me} & \texttt{ns1/*/foo}         & --       & no  \\
    \texttt{\textasciitilde/*/foo/} & \texttt{ns1} & \texttt{me} & \texttt{me/*/foo}    & --       & no  \\
    \texttt{/\textasciitilde/*/foo/} & \texttt{ns1} & \texttt{me} & \texttt{\textasciitilde/*/foo} & -- & no \\
\end{CyphalSimpleTable}

Note in the last row that an absolute name beginning with \texttt{/}
short-circuits home expansion, so the tilde is preserved as a literal
character in the resolved name.

Next, a remap table is installed. A matched remap rule replaces the user-supplied
name with the rule's \texttt{to} string, which is then fed through the same
pipeline; a matched rule also discards any pin suffix the user may have
supplied. The pin on a rule's \texttt{to} string, if present, is honoured
instead.

\begin{CyphalSimpleTable}[wide]{Name resolution with remap}{|l l l l l r X|}
    Name                 & From           & To                & Namespace & Home    & Pin    & Remark \\
    \texttt{foo/bar}     & \texttt{foo/bar} & \texttt{zoo}     & \texttt{ns} & \texttt{me} & --   & Relative remap; resolved name is \texttt{ns/zoo}. \\
    \texttt{foo/bar}     & \texttt{foo/bar} & \texttt{zoo\#123} & \texttt{ns} & \texttt{me} & $123$ & Pinned relative remap; resolved name is \texttt{ns/zoo}. \\
    \texttt{foo/bar\#456} & \texttt{foo/bar} & \texttt{zoo}     & \texttt{ns} & \texttt{me} & --   & User pin discarded by the matched rule; resolved name is \texttt{ns/zoo}. \\
    \texttt{foo/bar}     & \texttt{foo/bar} & \texttt{/zoo}    & \texttt{ns} & \texttt{me} & --   & Absolute remap; namespace is ignored; resolved name is \texttt{zoo}. \\
    \texttt{foo/bar}     & \texttt{foo/bar} & \texttt{\textasciitilde/zoo} & \texttt{ns} & \texttt{me} & -- & Homeful remap; the tilde expands to \texttt{me}; resolved name is \texttt{me/zoo}. \\
\end{CyphalSimpleTable}

The following inputs cause resolution to fail:

\begin{description}
    \item[\texttt{"foo bar\textbackslash{}nbaz"}] Contains a space (ASCII $32$)
    and a control character (ASCII $10$); non-printable and space characters
    are not allowed.
    \item[\texttt{"foo/*/bar\#123"}] A pattern name cannot be pinned; the
    \texttt{\#} suffix is not permitted on names containing pattern tokens.
    \item[\texttt{""} (empty)] An empty name is not allowed.
    \item[\texttt{"\#1234"}] After the pin expression is stripped, the
    remaining name is empty, which is not allowed.
    \item[A name longer than 200 bytes after normalization] Rejected regardless
    of how much buffer space is available.
\end{description}

\section{Topic identity and subject-ID mapping}\label{sec:session_identity}

Every topic has a deterministic 64-bit \emph{topic hash} derived from its normalized
resolved name and a transport-visible \emph{subject-ID} derived from the hash and
its eviction counter.
The mapping is purely arithmetic and uses no shared state other than the parameters
normalized in this section; every node computes the same value from the same inputs.

\subsection{The topic hash function}\label{sec:session_rapidhash}

The topic hash is computed as $\mathrm{rapidhash}(\mathit{name},\ \mathit{seed}=0)$,
where \textit{name} is the UTF-8 byte sequence of the normalized resolved name
(no NUL terminator) and $\mathrm{rapidhash}$ is the 64-bit
\textbf{rapidhash V3} function\footnote{%
    rapidhash is a fast, high-quality, platform-independent hash function
    developed by Nicolas De Carli and distributed under the MIT license at
    \href{https://github.com/Nicoshev/rapidhash}{github.com/Nicoshev/rapidhash}.
    It is derived from the wyhash algorithm by Wang Yi.
    Cyphal v1.1 pins the V3 variant of rapidhash as its normative hash function;
    the Python transcription below has been verified bit-identical to the
    reference C implementation for every test vector in this section.
} with the default secret specified below.

The default secret consists of eight 64-bit constants:

\begin{CyphalSimpleTable}{Default secret for the rapidhash V3 topic hash function}{|l l|}
    Index   & Value (hexadecimal)      \\
    0       & \texttt{0x2d358dccaa6c78a5} \\
    1       & \texttt{0x8bb84b93962eacc9} \\
    2       & \texttt{0x4b33a62ed433d4a3} \\
    3       & \texttt{0x4d5a2da51de1aa47} \\
    4       & \texttt{0xa0761d6478bd642f} \\
    5       & \texttt{0xe7037ed1a0b428db} \\
    6       & \texttt{0x90ed1765281c388c} \\
    7       & \texttt{0xaaaaaaaaaaaaaaaa} \\
\end{CyphalSimpleTable}

Multi-byte integers within the hash function are read from the input in
\emph{little-endian} byte order on every platform; big-endian implementations shall
emulate this ordering. Wherever the resulting 64-bit hash integer appears in a
session-layer header, it is serialized as little-endian.

\begin{remark}
    Multiple rapidhash variants exist in the wild
    (V1, V2, V3, V3-micro, V3-nano, wyhash, \ldots) and they are not wire-compatible
    with one another. Cyphal v1.1 pins down the full V3 variant only, with the
    default seed and the default secret listed above.
\end{remark}

A normative Python transcription of the reference C implementation is given in
appendix~\ref{sec:appendix_rapidhash}. It is bit-identical to the reference
implementation on every platform and may be used directly to reproduce the
test vectors below.

The following test vectors are authoritative. An implementation that fails to
reproduce any of them is not interoperable with the rest of the Cyphal network and
cannot participate in the topic allocation protocol.

\begin{CyphalSimpleTable}{Rapidhash V3 test vectors}{|l r l|}
    Input (ASCII)                            & Length & \texttt{rapidhash(input)}  \\
    \texttt{""} (empty)                      & 0      & \texttt{0x0338dc4be2cecdae} \\
    \texttt{"a"}                             & 1      & \texttt{0x599f47df33a2e1eb} \\
    \texttt{"abc"}                           & 3      & \texttt{0xcb475beafa9c0da2} \\
    \texttt{"topic"}                         & 5      & \texttt{0x0164a16804e84989} \\
    \texttt{"sensors/imu"}                   & 11     & \texttt{0xbb402fd6c693c60b} \\
    \texttt{"sensors/engine/temperature"}    & 26     & \texttt{0x32d2c4bd5c5485e7} \\
    \texttt{"head/camera/upper\_left"}       & 22     & \texttt{0x28f52efde90dec0a} \\
    \texttt{"system/primary/sensor/0/data"}  & 28     & \texttt{0x57627b4b782b47af} \\
    \texttt{"ins/0/data"}                    & 10     & \texttt{0xdf3339ab99373b98} \\
    \texttt{"1234"}                          & 4      & \texttt{0xa5ae810dbac9be9b} \\
    \texttt{"65535"}                         & 5      & \texttt{0x2c5a542e44fa7c20} \\
    200 $\times$ \texttt{'x'}                & 200    & \texttt{0x25adb114ca9ea76a} \\
\end{CyphalSimpleTable}

\subsection{Subject-ID mapping}\label{sec:session_subject_id_mapping}

The subject-ID assigned to a non-pinned topic is a function of its hash $h$, its
eviction counter $e$, and the transport's subject-ID modulus $M$:
\begin{equation*}
    \mathit{subject\_id}(h, e) = 8192 + ((h + e \cdot e) \bmod M)
\end{equation*}
All arithmetic is performed on 64-bit unsigned integers with wraparound;
implementations using signed types shall take care not to sign-extend the hash.
The modulus $M$ is a prime satisfying $M \bmod 4 = 3$, chosen per
section~\ref{sec:session_modulus}; its selection ensures that the
$\{e^2 \bmod M : e \in \mathbb{N}\}$ sequence is a valid quadratic-probing
sequence.

A \emph{pinned} topic is one whose eviction counter lies in the reserved range
$[\text{UINT32\_MAX}-8191,\ \text{UINT32\_MAX}]$.
Its subject-ID is obtained by $\mathit{subject\_id} = \text{UINT32\_MAX} - e$,
which lies in the pinned range $[0, 8191]$.
Pinned topics otherwise behave as ordinary CRDT records and participate in
arbitration on equal terms with non-pinned topics.

\begin{remark}
    The high-eviction encoding of pinned topics serves a second purpose: it acts as a
    wrap-around guard for the eviction counter itself.
    Under any pathological churn that drove a non-pinned topic's eviction counter
    past $2^{32} - 8192$, the topic would enter the pinned range and be arrested at
    subject-ID 8191 rather than overflowing the counter and disrupting consensus.
    The 32-bit counter is wide enough that such growth is not reachable in practice:
    at a sustained one eviction per second it takes over 136 years to overflow.
\end{remark}

\subsection{Subject-ID ranges}\label{sec:session_subject_id_ranges}

The session layer partitions the entire subject-ID space exposed by the
transport into four disjoint ranges, described in the table below.
All values refer to the numerical subject-ID carried in the transport header;
the topic $\to$ subject-ID mapping is defined in
section~\ref{sec:session_subject_id_mapping}.

\begin{CyphalSimpleTable}{Subject-ID space partition}{|l X|}
    Range                          & Usage \\
    $[0,\ 8191]$                   & Pinned topics. Retained as the primary path for interoperation with Cyphal/CAN v1.0 nodes, which only understand this range. \\
    $[8192,\ 8191{+}M]$            & Auto-allocated topics. $M$ is a transport-specific prime modulus satisfying $M \bmod 4 = 3$, defined in section~\ref{sec:session_modulus}. \\
    $(8191{+}M,\ B)$               & Gossip shard subjects used for background CRDT dissemination (section~\ref{sec:session_gossip}). \\
    $B$                            & The \emph{broadcast subject}, used for broadcast gossip and for scout service discovery (section~\ref{sec:session_scout}). \\
\end{CyphalSimpleTable}

where $B = 2^{\lceil\log_2(8191 + M + 1)\rceil} - 1$ is the highest subject-ID
representable in the transport's subject-ID width.
Every node is expected to subscribe to the broadcast subject; the transport's
implicit participant discovery mechanism (section~\ref{sec:transport}) is a
side effect of that universal subscription.

On the broadcast subject and on any gossip shard subject, the transport layer
is permitted to relax two of the requirements stated in
section~\ref{sec:transport}: deduplication of incoming transfers is not
mandatory, and support for transfers exceeding a single transport frame is not
required. The session layer tolerates occasional duplication on these
subjects, and constrains its own protocol messages on them to fit within a
single transport frame.

\begin{remark}
    The relaxation is important for scalability on resource-constrained
    transports. Deduplication state on the broadcast subject would otherwise
    grow linearly with the number of peers, which is hostile to small nodes;
    multi-frame reassembly buffers on the gossip shards would impose a
    per-shard memory cost that a large network cannot afford.
\end{remark}

\subsection{Subject-ID modulus}\label{sec:session_modulus}

The modulus $M$ depends on the width of the subject-ID field supported by the
transport. It shall be prime, satisfy $M \bmod 4 = 3$, and be the largest such
prime not exceeding $(\text{transport maximum subject-ID}) - 8191$.
The following values are normative:

\begin{CyphalSimpleTable}{Subject-ID moduli for supported transport widths}{|l l r X|}
    Name    & Width   & Modulus $M$     & Notes                                                             \\
    16-bit  & 16 bits & $57203$         & Usable on all Cyphal transports, including Cyphal/CAN.             \\
    23-bit  & 23 bits & $8378431$       & Usable on Cyphal/UDP; not compatible with Cyphal/CAN.              \\
    32-bit  & 32 bits & $4294954663$    & Reserved for future 32-bit transports; not compatible with the above. \\
\end{CyphalSimpleTable}

All nodes on the same logical Cyphal network shall use the same modulus.
In a heterogeneously redundant deployment where different transports are used in
parallel, the smallest applicable modulus shall be chosen so that every subject-ID
is representable on every member transport.

For the 16-bit modulus the derived subject-ID ranges are $[0, 8191]$ for pinned
topics, $[8192, 65394]$ for auto-allocated topics, $[65395, 65534]$ for the
140 gossip shards, and $65535$ for the broadcast subject.
For the 23-bit modulus the ranges are $[0, 8191]$, $[8192, 8386622]$,
$[8386623, 8388606]$ ($1984$ gossip shards), and $8388607$ respectively.

The following test vectors lock down the mapping formula for typical inputs.
The evictions column traces the first three entries of the quadratic-probing
sequence $(e = 0, 1, 2)$; subsequent entries step by $5, 7, 9, \ldots$.

\begin{CyphalSimpleTable}{Subject-ID test vectors}{|l l r r r|}
    Name               & Hash                 & $e$ & $M=57203$ & $M=8378431$ \\
    \texttt{""}        & \texttt{0x0338dc4be2cecdae} & 0 & $21600$  & $6520572$ \\
    \texttt{""}        & \texttt{0x0338dc4be2cecdae} & 1 & $21601$  & $6520573$ \\
    \texttt{""}        & \texttt{0x0338dc4be2cecdae} & 2 & $21604$  & $6520576$ \\
    \texttt{"topic"}   & \texttt{0x0164a16804e84989} & 0 & $38082$  & $1049799$ \\
    \texttt{"topic"}   & \texttt{0x0164a16804e84989} & 1 & $38083$  & $1049800$ \\
    \texttt{"topic"}   & \texttt{0x0164a16804e84989} & 2 & $38086$  & $1049803$ \\
    \texttt{"sensors/imu"} & \texttt{0xbb402fd6c693c60b} & 0 & $54214$ & $8263034$ \\
    \texttt{"sensors/imu"} & \texttt{0xbb402fd6c693c60b} & 1 & $54215$ & $8263035$ \\
    \texttt{"sensors/imu"} & \texttt{0xbb402fd6c693c60b} & 2 & $54218$ & $8263038$ \\
    \texttt{"ins/0/data"} & \texttt{0xdf3339ab99373b98} & 0 & $28748$ & $8266023$ \\
    \texttt{"ins/0/data"} & \texttt{0xdf3339ab99373b98} & 1 & $28749$ & $8266024$ \\
    \texttt{"ins/0/data"} & \texttt{0xdf3339ab99373b98} & 2 & $28752$ & $8266027$ \\
\end{CyphalSimpleTable}

\section{Topic allocation CRDT}\label{sec:session_crdt}

The session layer maintains the topic-to-subject-ID mapping through a conflict-free
replicated data type (CRDT) shared among all nodes by periodic gossip.
Under the assumptions listed below, every node's local copy of the mapping converges
to the same collision-free, non-divergent configuration within a bounded time after
the last state-changing event.

This section is normative for the arbitration semantics of the mapping.
The primary source of truth is the TLA$^+$ model in \texttt{model/tlaplus/Core.tla}
of the Cy reference implementation; the operators and theorems described here are
a faithful transcription of that model, and the C reference implementation has been
shown to match it under the stated assumptions.

\subsection{Topic record}

Each topic the node is aware of is represented by a triple
\[
    t = (\mathit{hash},\ \mathit{evictions},\ \mathit{age})
\]
where:

\begin{description}
    \item[\textit{hash}] is the 64-bit topic hash defined in
    section~\ref{sec:session_rapidhash}.
    \item[\textit{evictions}] is a 32-bit unsigned integer incremented each time the
    topic loses arbitration on its current candidate subject-ID. It serves as a
    Lamport clock. The reserved range at the top of the counter encodes pinned
    topics as described in section~\ref{sec:session_subject_id_mapping}.
    \item[\textit{age}] is a local real or logical clock counting roughly seconds
    since the topic was first seen on the network. It is not transmitted directly
    on the wire; gossip instead carries its base-2 logarithm (see below).
\end{description}

The subject-ID of a topic is the pure function
$\mathit{subject\_id}(\mathit{hash}, \mathit{evictions})$ defined in
section~\ref{sec:session_subject_id_mapping}.

\subsection{Log-age}

CRDT operators use the quantity
\[
    \mathit{lage}(t) = \lfloor \log_2(\mathit{age\_seconds}(t)) \rfloor,
    \qquad \mathit{lage}(t) = -1 \text{ when } \mathit{age\_seconds}(t) = 0
\]
in place of the raw age.
Taking the floor of the logarithm slows the arbitration-relevant clock as topics
grow older, so that the CRDT is stable on long timescales even though the underlying
age counter never stops advancing. The encoded range is $[-1,\ 35]$, where
$\mathit{lage} = 35$ corresponds to approximately $2^{35}$ seconds, or over a
thousand years.

\begin{remark}
    The naive approach of using the raw age in arbitration would make any two
    quiescent nodes disagree on priorities every time a clock edge crosses an
    integer second boundary, which breaks the CRDT's eventual consistency.
    Reducing the comparison quantity to $\lfloor \log_2(\mathit{age}) \rfloor$
    makes the set of points where priority can change sparse on quiescent states,
    satisfying the stability assumption A9 below.
\end{remark}

Gossip messages carry a topic's $\mathit{lage}$ value rather than its age.
On reception, a node merges the incoming $\mathit{lage}$ into its local age by
adjusting the local age field so that
\[
    \mathit{local\_age} \leftarrow \max\left(\mathit{local\_age},\ 2^{\mathit{remote\_lage}}\right).
\]
This is a monotonic, idempotent join: repeated application is a no-op, and the
result is independent of the order in which gossip messages are processed.

\subsection{Arbitration}

Two topics are in \emph{collision} when they have distinct hashes but map to the
same subject-ID. Two records with the same hash but different eviction counters
are in \emph{divergence}; these are different opinions about the same logical topic
that must be reconciled.

\begin{description}
    \item[Collision arbitration] Older $\mathit{lage}$ wins. On a tie, the record
    with the numerically lower hash wins.
    \item[Divergence arbitration] Older $\mathit{lage}$ wins. On a tie, the record
    with the higher eviction count wins, because a higher counter means the record
    has observed more prior conflicts and is therefore more up to date.
\end{description}

These two predicates are defined as follows, verbatim from
\texttt{model/tlaplus/Core.tla}. \texttt{Log2Floor(x)} is
$\lfloor \log_2 x \rfloor$ with the special case $\mathrm{Log2Floor}(0) = -1$;
\texttt{\#} is the TLA\textsuperscript{+} inequality operator, $\mathtt{=}$ is
equality, $\mathtt{/\char`\\}$ is logical AND, and $\mathtt{\char`\\/}$ is
logical OR.

\begin{minted}{text}
LeftWinsCollision(a, b) == IF   Log2Floor(a.age) # Log2Floor(b.age)
                           THEN Log2Floor(a.age) > Log2Floor(b.age)
                           ELSE a.hash < b.hash

LeftWinsDivergence(a, b) == \/  Log2Floor(a.age) > Log2Floor(b.age)
                            \/ (Log2Floor(a.age) = Log2Floor(b.age) /\ a.evictions > b.evictions)
\end{minted}

\subsection{AllocateTopic operator}

\texttt{AllocateTopic} inserts a topic into the node's local topic set, resolving
any resulting collision by cascading evictions through the hash space.
The operator is recursive; termination follows from the bounded-escape property in
assumption A3 below. The TLA\textsuperscript{+} definition below is verbatim
from \texttt{model/tlaplus/Core.tla} and is reproduced in full in
appendix~\ref{sec:appendix_core_tla}.

\begin{minted}{text}
RECURSIVE AllocateTopic(_, _)
AllocateTopic(t, topics) ==
    LET ts == RemoveTopic(t.hash, topics)
        x == GetBySubjectID(SubjectID(t.hash, t.evictions), ts)
        Evicted(z) == [hash |-> z.hash, evictions |-> 1 + z.evictions, age |-> z.age]
    IN   IF x = Nothing             THEN {t} \cup ts
    ELSE IF LeftWinsCollision(t, x) THEN AllocateTopic(Evicted(x), {t} \cup ts)
    ELSE                                 AllocateTopic(Evicted(t), ts)
\end{minted}

The operator is invoked whenever a new local topic is created, whenever a gossip
forces the local record for a topic to adopt new eviction values, and whenever
a cascade from another topic displaces the current record.

\subsection{AcceptGossip operator}

\texttt{AcceptGossip} applies a received gossip message to the node's local topic
set. The operator distinguishes two cases: the gossiped topic is known locally
(its hash matches a local record) or it is not. The TLA\textsuperscript{+}
definitions below are verbatim from \texttt{model/tlaplus/Core.tla}.
\texttt{Nothing} denotes the absence of a matching record;
\texttt{FloorToPow2(x)} returns $2^{\lfloor \log_2 x \rfloor}$ with the
special case $\mathrm{FloorToPow2}(0) = 0$.

\begin{minted}{text}
AcceptGossip(remote, topics) ==
    IF GetByHash(remote.hash, topics) # Nothing     \* If the topic exists locally
    THEN AcceptGossip_Divergence(remote, topics)    \* Divergence resolver ensures collision-freedom
    ELSE AcceptGossip_Collision(remote, topics)     \* We don't know this topic, just ensure collision-freedom

AcceptGossip_Divergence(remote, topics) ==
    LET hash == remote.hash
        local == GetByHash(hash, topics)
        \* The serialization protocol floors the remote age to pow2 before transmission.
        new_age == Max({IF local # Nothing THEN local.age ELSE 0, FloorToPow2(remote.age)})
    IN
    IF local # Nothing THEN
        IF local.evictions # remote.evictions /\ LeftWinsDivergence(remote, local)
        THEN AllocateTopic([hash|->hash, evictions|->remote.evictions, age|->new_age], topics)
        ELSE              {[hash|->hash, evictions|-> local.evictions, age|->new_age]} \cup RemoveTopic(hash, topics)
    ELSE topics

AcceptGossip_Collision(remote, topics) ==
    LET local == GetBySubjectID(SubjectID(remote.hash, remote.evictions), topics)
    IN IF local # Nothing /\ LeftWinsCollision(remote, local)
    THEN AllocateTopic([hash |-> local.hash, evictions |-> local.evictions + 1, age |-> local.age], topics)
    ELSE topics
\end{minted}

\begin{remark}[breakable]
    A naive implementation could always run the collision resolver after the
    divergence resolver, i.e.\ compute
    \texttt{AcceptGossip\_Collision(remote, AcceptGossip\_Divergence(remote, topics))}.
    This is correct but suboptimal: the local record may already win the
    divergence arbitration against the remote, in which case the remote record
    should be ignored entirely; running the collision resolver afterwards can make
    a local record lose a collision to a stale remote record and trigger
    unnecessary eviction-counter increments. The divergence-first rule above
    avoids that outcome -- when the local record wins the divergence arbitration,
    the collision pass is skipped because divergence resolution implicitly keeps
    the local mapping collision-free.
\end{remark}

\subsection{Pinned topics}

Pinned topics are ordinary CRDT records with their eviction counter pre-set to
$\text{UINT32\_MAX} - \text{pinned\_subject\_id}$, as defined in
section~\ref{sec:session_subject_id_mapping}.
They participate in arbitration on the same terms as non-pinned topics: an older
non-pinned topic can displace a younger pinned newcomer, and an older pinned topic
can displace a younger conflicting pinned topic.
Pinning is therefore a \emph{request}, not a guarantee.

\subsection{Assumptions}

The convergence guarantee of the allocation protocol rests on the following
assumptions. Some are structural properties of the mapping function and are
satisfied by construction; others are deployment constraints that compliant
implementations and system integrators are required to uphold. The numbering is
consistent with the accompanying proof in \texttt{model/proof.md}.

\begin{description}
    \item[A1. Common deterministic mapping] All nodes use the same subject-ID
    mapping function and the same parameters -- in particular, the same modulus.
    \item[A2. Hash identity] Active topic names are hash-unique. The probability of
    a 64-bit hash collision is negligible and is ignored by the protocol.
    \item[A3. Bounded escape] For any hash, any eviction value, and any blocked
    subject-ID set smaller than the capacity bound $K$, there exists a small
    increment within $K$ steps that places the topic at a subject-ID outside the
    blocked set. For quadratic probing on a prime modulus $M$ with $M \bmod 4 = 3$
    this holds with $K \approx \lfloor M / 2 \rfloor + 1$.
    \item[A4. Capacity bound] The number of active non-pinned topics $N$ satisfies
    $N \le K$.
    \item[A5. Gossip fairness] Every live topic record is eventually received and
    processed by every relevant node infinitely often.
    \item[A6. Progress fairness] Each node keeps executing the protocol; resource
    exhaustion or scheduling starvation does not halt a live node.
    \item[A7. Resource sufficiency] Each node has enough memory and CPU to process
    and complete the local cascade triggered by any single allocation step.
    \item[A8. Eviction counter safety] Eviction counters do not overflow in
    reachable executions.
    \item[A9. Stability window] Stable windows longer than the worst-case repair
    time $T_{\mathrm{repair}}$ exist, during which the active topic set and
    pairwise arbitration priorities do not flip.
\end{description}

\subsection{Convergence}

Under assumptions A1--A9, any execution restricted to a single stability window
converges in finite time to a state that is free of collisions and divergences.
The argument is given in full in \texttt{model/proof.md}; the key idea is that
higher-priority topics are never displaced by lower-priority topics, so the
arbitration induces a fixed total order during the window and each topic
stabilizes by induction on its priority rank.

When the stability window ends -- for example, because a long-lived topic's
$\mathit{lage}$ steps up and flips a tie -- some transient churn may briefly
reappear. Assumption A9 guarantees that the next stable window is long enough to
repair the churn before the next flip, and the overall execution remains eventually
consistent.

\subsection{Capacity}

For a quadratic probing sequence on a prime modulus $M$ with $M \bmod 4 = 3$ the
collision-free probing capacity is approximately $\lfloor M / 2 \rfloor$
(assumption A3). The number of active non-pinned topics should not exceed this
bound; above it, the arbitration behaviour remains well-defined but convergence
within a single stability window is no longer guaranteed.

\section{Gossip dissemination}\label{sec:session_gossip}

Nodes keep their local CRDT states (section~\ref{sec:session_crdt}) aligned with
each other by exchanging \emph{gossip} messages.
A gossip carries a single topic record and is sent over one of four carriers
described below.

\subsection{Gossip carriers}

\begin{description}
    \item[Broadcast gossip] Multicast on the broadcast subject
    (section~\ref{sec:session_subject_id_mapping}). Every node is expected to be
    attached to this subject, so a broadcast gossip is eventually visible to every
    participant; broadcast traffic is the primary mechanism for escaping
    partitioned or newly-joined subsets. The broadcast subject is also used for
    scout messages (section~\ref{sec:session_scout}).

    \item[Sharded gossip] Multicast on a \emph{gossip shard} subject-ID.
    A topic is assigned to the shard
    \[
        \mathit{shard\_index} = \mathit{hash} \bmod \mathit{gossip\_shard\_count}
    \]
    where $\mathit{gossip\_shard\_count}$ is the number of subject-IDs between
    the auto-allocated range and the broadcast subject (see
    section~\ref{sec:session_subject_id_mapping}). A node joins every gossip shard
    that any of its local topics maps into, and no others.

    \item[Inline gossip] A gossip state carried in the header of a
    \texttt{msg\_be} or \texttt{msg\_rel} application message
    (section~\ref{sec:session_headers}). Inline gossip is observed by every
    subscriber of the topic and converges the CRDT state for that topic without
    spending a separate message on it. Handling of inline gossip is recommended
    but not required for correctness.

    \item[Unicast gossip] A single gossip sent to a specific remote node,
    used for scout responses and for targeted repair of a known divergence.
\end{description}

\subsection{Schedules}

A node that holds at least one explicit topic emits gossip on three schedules:
periodic, urgent, and inline. Implicit topics
(section~\ref{sec:session_patterns}) do not emit gossip.

\paragraph{Periodic gossip} Every explicit topic emits its current state at a
randomized interval around a nominal period (recommended default: $5$ seconds,
randomized by $\pm$ one eighth).
Every $N$-th emission goes on the broadcast subject rather than the topic's
gossip shard, so the broadcast lane sees a slow but guaranteed sweep of all
topics on the network (recommended $N = 10$).
Between broadcast sweeps the topic's shard is used, which scales better on large
networks because each node only joins a fraction of the shards.

\paragraph{Urgent gossip} Whenever a topic's local subject-ID occupancy changes
-- a new topic is created, an incoming gossip forces a re-allocation, or a
cascade evicts a record -- a broadcast gossip is scheduled with a small random
delay $W_u$ (section~\ref{sec:session_gossip_delays}). If an equivalent or
dominating gossip from another node arrives before the timer fires, the pending
urgent gossip is suppressed.

\paragraph{Inline gossip} A node that publishes a \texttt{msg\_be} or
\texttt{msg\_rel} message piggybacks the current gossip state of the message's
topic in the header. A receiver that elects to process inline gossip applies it
to its local CRDT in the usual way; this is the fastest possible CRDT update on
active topics because it consumes no additional network capacity.

\subsection{Suppression rules}\label{sec:session_gossip_suppression}

Gossip schedules are designed so that the network-wide rate of gossip traffic on
any given topic stays roughly constant regardless of the number of nodes holding
it. To that end, nodes apply two suppression rules when incoming gossip indicates
that some other node is already covering the same information:

\begin{itemize}
    \item \textbf{Cancel a pending urgent gossip} when an incoming gossip arrives
    on the broadcast subject whose arbitration priority dominates the state the
    urgent gossip was intended to publish. The incoming broadcast already
    carries the information forward, so repeating it would waste network
    capacity.

    \item \textbf{Slow down periodic gossip} when an incoming non-inline gossip
    matches the local $\mathit{lage}$ and eviction counter for the same topic.
    Another node is already covering the topic at the current cadence; the
    local node can afford to reschedule its next emission with a longer
    interval.
\end{itemize}

\begin{remark}
    Inline gossip is deliberately excluded from the periodic slowdown rule.
    Because inline gossip is piggybacked on application messages, its rate is
    determined by the application's publication rate, not by the gossip
    scheduler; allowing it to suppress periodic gossip would couple the CRDT
    repair rate to the application traffic pattern.
\end{remark}

\subsection{Randomized delay windows}\label{sec:session_gossip_delays}

The urgent window $W_u$ and the startup window $W_s$ are both tunable. The
recommended orders of magnitude below are derived from the analytic model in
\texttt{model/gossip\_propagation.ipynb} of the reference implementation.

Let $\delta_{\max}$ be the worst-case propagation and processing delay for a
gossip between two nodes on the network.

\begin{description}
    \item[Urgent window $W_u$] Controls the amount of jitter inserted before an
    urgent gossip. If $m_u$ nodes independently schedule an urgent repair for
    the same event, the expected number of redundant urgent emissions is
    approximately $(m_u - 1) \cdot \delta_{\max} / W_u$.
    In typical deployments $m_u$ is small because a given collision or
    divergence is observed by only a few nodes; setting
    $W_u \approx \delta_{\max}$ is usually sufficient.

    \item[Startup window $W_s$] Controls the amount of jitter inserted when a
    node comes online and schedules its first round of gossip.
    Wider windows reduce the probability of simultaneous emissions from many
    startups but saturate quickly; recommended range is
    $[\delta_{\max},\ 10 \cdot \delta_{\max}]$.
\end{description}

On a typical local network with $\delta_{\max}$ on the order of tens of
milliseconds, choosing $W_u$ on the order of tens of milliseconds and $W_s$ on
the order of tens to a few hundred milliseconds is sufficient.

\subsection{Implicit topic retirement}

An \emph{implicit} topic is one that a node holds only because a local pattern
subscription (section~\ref{sec:session_patterns}) matched a remote gossip. An
implicit topic does not emit gossip, and it is automatically retired once it has
observed no messages and no gossip activity for the implicit-topic retention
window. The recommended retention window is ten minutes; it should be much
longer than the worst-case periodic gossip interval so that transient silence
on an active topic does not cause premature retirement.

\section{Scout}\label{sec:session_scout}

The scout mechanism lets a node pull in topics that match a pattern without
waiting for periodic gossip rotation. A scout is a session-layer message whose
payload is a pattern string (section~\ref{sec:session_names}).

A scout message is broadcast on the broadcast subject using the
\texttt{scout} header (section~\ref{sec:session_headers}).
Every node that receives the message checks its local topics against the
pattern; for each match, the receiver replies with a \texttt{gossip} header
describing the matched topic. Replies are typically unicast back to the
originating node, using the unicast context provided by the transport with the
incoming scout, but broadcasting is also permitted if that is easier for the
responder.

A scout with an empty pattern is reserved for a future extension -- a request
for a network-wide health summary -- and shall be discarded by current
implementations.

\section{Session headers and framing}\label{sec:session_headers}

Every session-layer message begins with a fixed-size header of 24 bytes,
followed by a type-specific payload. Multi-byte integer fields are encoded in
little-endian byte order and are positioned to favour natural alignment on
four- and eight-byte boundaries. The wire layouts below are given in DSDL
notation.
Fields marked \emph{incompatibility} shall be transmitted as
zero, and any non-zero value on reception indicates a protocol revision the
receiver does not understand: the receiver shall discard the message.

\subsection{Header types}

The first byte of every header is the header type. Types are allocated as
follows; values 10 through 255 are reserved.

{\footnotesize
\begin{CyphalSimpleTable}[wide]{Session-layer header types}{|r l l X|}
    Value & Name             & Direction                         & Purpose                                        \\
    0     & \texttt{msg\_be}    & publisher $\to$ subscriber(s), multicast or unicast & Best-effort message publication.                \\
    1     & \texttt{msg\_rel}   & publisher $\to$ subscriber(s), multicast or unicast & Reliable message publication, acknowledged.     \\
    2     & \texttt{msg\_ack}   & subscriber $\to$ publisher, unicast              & Positive acknowledgement of a reliable message. \\
    3     & \texttt{msg\_nack}  & subscriber $\to$ publisher, unicast              & Negative acknowledgement; no matching subscriber.\\
    4     & \texttt{rsp\_be}    & responder $\to$ publisher, unicast               & Best-effort response to a received message.     \\
    5     & \texttt{rsp\_rel}   & responder $\to$ publisher, unicast               & Reliable response, acknowledged.                \\
    6     & \texttt{rsp\_ack}   & publisher $\to$ responder, unicast               & Positive acknowledgement of a reliable response.\\
    7     & \texttt{rsp\_nack}  & publisher $\to$ responder, unicast               & Negative acknowledgement; application no longer listening. \\
    8     & \texttt{gossip}     & any, multicast or unicast                        & Topic allocation CRDT gossip.                   \\
    9     & \texttt{scout}      & any, broadcast                                   & Pattern service-discovery query.                \\
\end{CyphalSimpleTable}
}

\subsection{Message headers}\label{sec:session_header_msg}

The \texttt{msg\_be} and \texttt{msg\_rel} headers share the same layout. Each
message carries not only the application payload but also the sender's
current CRDT state for the topic, so that the message doubles as an inline
gossip for its own topic (section~\ref{sec:session_gossip}).

\begin{minted}{python}
uint8   type             # 0 for msg_be, 1 for msg_rel.
void8
uint8   incompatibility
int8    topic_log_age    # floor(log2(topic_age_seconds)) if topic_age_seconds > 0, else -1.
uint32  evictions        # For pinned topics, 0xFFFFFFFF - subject_id.
uint64  topic_hash       # Subject allocation collision detection and immediate consensus updates.
uint64  tag              # Random-initialized, wraps around. Used for acknowledgement
                         # correlation and, optionally, for ordering recovery per remote.
# Payload follows.
\end{minted}

\subsection{Message acknowledgement headers}\label{sec:session_header_msg_ack}

The \texttt{msg\_ack} and \texttt{msg\_nack} headers are sent in response to a
\texttt{msg\_rel}; they carry no payload. The priority level of the
acknowledgement shall match the priority level of the original message.

\begin{minted}{python}
uint8   type             # 2 for msg_ack, 3 for msg_nack.
void24
uint32  incompatibility
uint64  topic_hash       # From the acknowledged message.
uint64  tag              # From the acknowledged message.
\end{minted}

\subsection{Response headers}\label{sec:session_header_rsp}

The \texttt{rsp\_be} and \texttt{rsp\_rel} headers carry a response to a
message previously received by the responder. The \texttt{rsp\_ack} and
\texttt{rsp\_nack} headers acknowledge a reliable response and carry no
payload. All four share the same layout.

\begin{minted}{python}
uint8   type             # 4 for rsp_be, 5 for rsp_rel, 6 for rsp_ack, 7 for rsp_nack.
uint8   tag              # Chosen by the responder arbitrarily for acknowledgement correlation.
uint48  seqno            # Starts at zero for the first response to a given message and
                         # increments by one per successive response, enabling streaming.
uint64  topic_hash       # Topic hash of the published message this response pertains to.
uint64  message_tag      # Tag of the published message this response pertains to.
# Payload follows, except for rsp_ack and rsp_nack.
\end{minted}

\subsection{Gossip header}\label{sec:session_header_gossip}

The \texttt{gossip} header carries one CRDT record
(section~\ref{sec:session_crdt}) together with the topic's normalized
resolved name. The message has no additional payload.

\begin{minted}{python}
uint8   type             # 8.
void16
int8    topic_log_age    # floor(log2(topic_age_seconds)) if topic_age_seconds > 0, else -1.
uint32  incompatibility
uint64  topic_hash
uint32  topic_evictions  # For pinned topics, 0xFFFFFFFF - subject_id.
void24                   # Reserved, may extend the eviction counter in a future revision.
utf8[<=200] topic_name   # Length-prefixed (one byte prefix); normalized resolved name.
# Total wire size is 24 bytes plus the length of the topic name.
\end{minted}

\subsection{Scout header}\label{sec:session_header_scout}

The \texttt{scout} header carries a pattern string
(section~\ref{sec:session_names}). The receiver matches the pattern against
its local topics and, for each match, replies with a \texttt{gossip} message
(section~\ref{sec:session_scout}).

\begin{minted}{python}
uint8   type             # 9.
void24
uint32  incompatibility
uint64  incompatibility1
void56
utf8[<=200] pattern      # Length-prefixed (one byte prefix); pattern applied to normalized names.
# Total wire size is 24 bytes plus the length of the pattern.
\end{minted}

\subsection{Message tags}\label{sec:session_tags}

The 64-bit \texttt{tag} field carried in \texttt{msg\_be} and \texttt{msg\_rel}
is a per-publisher counter that is randomly initialized when a publisher is
created and is incremented by one for each message published on the topic.
The random initialization ensures that a node rebooting mid-operation does not
accidentally reuse tags that were emitted before the reboot; the counter is
permitted to wrap around.

The PRNG used for tag initialization shall have distinct state across quick
reboots. Nodes without a true random-number generator should seed a pseudorandom
generator by hashing the node's unique identifier with whatever entropy the
platform makes available -- for example, analog-to-digital converter noise,
real-time clock drift, or values stored in battery-backed memory.

Tags are not used for duplicate suppression on best-effort messages, which rely
on the transport layer's deduplication instead. They are used for duplicate
suppression and acknowledgement correlation on reliable messages
(section~\ref{sec:session_reliable}).

\section{Best-effort message exchange}\label{sec:session_best_effort}

A best-effort publication is a single \texttt{msg\_be} transfer. The publisher
performs the following steps atomically:

\begin{enumerate}
    \item Assign the next tag from the publisher's sequence.
    \item Serialize a \texttt{msg\_be} header whose inline CRDT fields carry
    the topic's current $\mathit{lage}$, eviction counter, and hash.
    \item Append the application payload.
    \item Submit the message to the transport for delivery on the topic's
    current subject-ID.
\end{enumerate}

The publisher retains no state after the submission returns.

On reception of a \texttt{msg\_be} transfer, a subscriber shall:

\begin{enumerate}
    \item Validate that the header's \texttt{(topic\_hash, evictions)} pair
    maps (via section~\ref{sec:session_subject_id_mapping}) to the subject-ID
    on which the message arrived. A mismatch indicates a malfunctioning
    sender; the message shall be discarded.
    \item Apply the inline gossip to the local CRDT state as described in
    section~\ref{sec:session_crdt}.
    \item Deliver the application payload to every matching local subscriber,
    if any.
\end{enumerate}

Best-effort messages do not require receive-side deduplication at the session
layer, because the transport already deduplicates transfers and the publisher
never retransmits.

\section{Reliable message delivery}\label{sec:session_reliable}

A reliable publication is carried by a \texttt{msg\_rel} transfer and is
acknowledged by every intended recipient with a \texttt{msg\_ack}. The
publisher retransmits until every known recipient has acknowledged or the
deadline expires.

\subsection{Associations}

A publisher tracks its active recipients in an \emph{association set}. Each
association records:

\begin{description}
    \item[\textit{remote\_id}] The network-wide identifier of the remote node.
    \item[\textit{unicast\_ctx}] The transport-specific metadata needed to send
    a unicast message back to the remote, as supplied by the transport layer
    with each incoming acknowledgement.
    \item[\textit{last\_seen}] Timestamp of the most recent acknowledgement
    received from the remote.
    \item[\textit{slack}] A miss counter incremented each time the remote
    fails to acknowledge a reliable message within the deadline. The
    association is forgotten when \textit{slack} exceeds a small limit
    (recommended default: $2$).
\end{description}

Associations are created opportunistically on receipt of a positive
acknowledgement for a reliable message, and are removed when they age out of
the session lifetime (section~\ref{sec:session_parameters}) or when their
slack counter exceeds the limit.

\subsection{First publication}

When a publisher sends its first reliable message on a topic it may have no
associations at all. The first message is multicast on the topic's subject-ID
and the publisher treats the arrival of \emph{any} positive acknowledgement as
success; further acknowledgements arriving after the first continue to
populate the association set in the background.

\subsection{Retransmission policy}

The publisher retransmits a reliable message using an exponential backoff
schedule. The first retransmission is scheduled one \emph{ACK baseline
timeout} after the initial attempt; each subsequent retransmission doubles the
timeout. The final attempt is capped by the publication deadline supplied with
the publication. Implementations may scale the ACK baseline timeout by the
message's priority level, for example to give higher-priority publications
tighter timeouts.

Initial attempts are always multicast, since the publisher does not yet know
which recipients are live. As acknowledgements arrive and unacknowledged
recipients remain, the publisher may elect to unicast the remaining
retransmissions directly to the holdouts using the transport unicast metadata
recorded in the association entries. This multicast-to-unicast fallback is a
heuristic chosen to reduce network and subscriber load; it is not required for
interoperability.

\subsection{Acknowledgements}

A \texttt{msg\_ack} is sent by any subscriber that accepts a reliable message
for processing. The acknowledgement is unicast back to the publisher at the
priority of the original message, and references the message by its topic
hash and tag.

A \texttt{msg\_nack} is sent only in response to a reliable \emph{unicast}
message for which the receiver has no matching subscriber. It tells the
publisher to prune the association promptly so that further retransmissions
are not wasted on a recipient that cannot accept the topic. \emph{Multicast}
reliable messages do not generate NACKs, because subscription membership is
implicit in the transport's multicast join state: a subscriber that is no
longer interested simply stops participating in the multicast group.

\begin{remark}
    The asymmetry between unicast and multicast NACKs is intentional. For
    multicast delivery, the publisher learns that a remote has left by
    observing its absence from future acknowledgements; for unicast delivery,
    no such implicit signal exists, so an explicit NACK is required.
\end{remark}

\subsection{Receive-side deduplication}

A reliable message may be delivered more than once by the transport if an
acknowledgement is lost and the publisher retransmits the same tag. The
receiver eliminates such duplicates by keeping a per-remote cache keyed by
\texttt{(remote\_id, tag)} and discarding any tag that has already been
accepted. Dedup state is soft and is aged out after the session lifetime
(section~\ref{sec:session_parameters}).

\subsection{Observable states of a reliable publication}

A reliable publication transitions through the following observable states,
which are named here to describe protocol behaviour rather than any API:

\begin{description}
    \item[\textsc{pending}] The initial state; set when the first attempt is
    submitted to the transport. The publication remains pending until a
    terminal event occurs.
    \item[\textsc{done, success}] Reached when at least one positive
    acknowledgement has been received and all remaining associations have
    either acknowledged or aged out of the publication's deadline.
    \item[\textsc{done, delivery failure}] Reached when the publication
    deadline expires without any positive acknowledgement.
    \item[\textsc{done, explicit rejection}] Reached on receipt of a
    \texttt{msg\_nack} for a unicast reliable message.
    \item[\textsc{done, transient lag}] Reached when the publication
    deadline is about to expire but the publisher cannot determine whether
    the retransmission schedule actually delivered the message due to local
    scheduler pressure. This state is informational: the application is
    warned that reachability information for this publication is
    incomplete.
\end{description}

\section{Deduplication and ordering}\label{sec:session_dedup}

Session-layer deduplication applies to reliable messages only, as described in
section~\ref{sec:session_reliable}. Best-effort messages are deduplicated by
the transport and no additional session-layer work is required.

A subscriber may optionally enable \emph{per-remote ordered delivery}.
When enabled, the subscriber uses the monotonic per-publisher \texttt{tag}
field to reconstruct the publication order within a bounded reorder window.
Messages whose tag lies within the window are delivered in increasing tag
order; messages whose tag lies outside the window because they arrived too
late are dropped. Ordering is guaranteed per remote publisher only, not across
remotes, because there is no global clock.

\section{RPC and streaming}\label{sec:session_rpc}

The session layer builds remote procedure calls on top of ordinary pub-sub
messages rather than introducing a separate service-ID abstraction. An RPC
request is a normal reliable publication whose recipients may emit one or more
unicast \emph{responses} back to the publisher. The number of responses per
request is unrestricted; streaming responses are therefore a zero-cost
extension of single-response RPC.

\subsection{Breadcrumbs}

When a subscriber receives a message, the transport supplies the unicast
metadata needed to reach the originating node. The subscriber captures this
information in a \emph{breadcrumb} consisting of:

\begin{itemize}
    \item \textit{remote\_id}: the network identifier of the original
    publisher;
    \item \textit{topic\_hash}: the hash of the topic the original message
    was delivered on;
    \item \textit{message\_tag}: the tag of the original message;
    \item \textit{seqno}: a 48-bit counter, initialized to zero, that is
    incremented by the subscriber once for each response it sends in reply
    to this message;
    \item \textit{unicast\_ctx}: the transport's unicast metadata for the
    original publisher.
\end{itemize}

A breadcrumb is sufficient to emit any number of responses to the original
publisher at a later time, without keeping the original message alive.

\subsection{Responses}

A response is carried by a \texttt{rsp\_be} or \texttt{rsp\_rel} transfer
(section~\ref{sec:session_header_rsp}). The response header carries:

\begin{itemize}
    \item \textit{topic\_hash} and \textit{message\_tag}, copied from the
    breadcrumb, which together identify the request;
    \item \textit{seqno}, the 48-bit response sequence number;
    \item \textit{response\_tag}, an arbitrary byte chosen by the responder
    for acknowledgement correlation (reliable responses only).
\end{itemize}

Streaming is achieved simply by sending multiple responses to the same
request, each with an incremented \textit{seqno}. The publisher observes a
monotonic stream of responses per remote responder.

\subsection{Reliable responses}

A responder that wishes to detect whether the original publisher is still
listening shall send its responses reliably (\texttt{rsp\_rel}). The
publisher acknowledges a reliable response with either a \texttt{rsp\_ack}
or a \texttt{rsp\_nack}:

\begin{description}
    \item[\texttt{rsp\_ack}] The publisher's application is still accepting
    responses for this request. The responder may continue to stream.

    \item[\texttt{rsp\_nack}] The publisher's application has closed the
    request and will discard any further responses. The responder should
    cease streaming promptly.
\end{description}

\subsection{Client-side deduplication and ordering}

The publisher tracks, for each remote responder, a small bitmap of recently
seen \textit{seqno} values and rejects duplicates within the bitmap window.
Responses whose \textit{seqno} falls outside the window because they are too
old are rejected.

Ordering is not reconstructed by default: a response with a larger
\textit{seqno} may be delivered before one with a smaller \textit{seqno}
on the same stream. An application that requires strict ordering within a
stream may implement a reorder buffer keyed on \textit{seqno}; pattern-based
subscribers may opt into such a buffer via the subscriber's ordering
configuration.

\subsection{Liveness monitoring}

The publisher may attach a \emph{liveness timeout} to a pending request.
If no response has been received from \emph{any} remote for longer than the
timeout, the publisher's observable state is moved into a \textsc{liveness}
warning that alerts the application and is automatically cleared when the
next response arrives. Liveness state is independent of the response queue:
a pending response that has not yet been consumed does not inhibit the
liveness warning, and clearing the warning does not drop the pending
response.

\subsection{Zombie requests}

After the publisher's application has closed a request, the session layer
may retain a small amount of state for the request briefly in order to
correctly handle duplicates of reliable responses that arrive after the
application has ceased to listen. A duplicate that arrives in this zombie
window is acknowledged with the same positive or negative acknowledgement
as the first time, without involving the application. The zombie window
expires with the session lifetime (section~\ref{sec:session_parameters}).

\begin{remark}
    The zombie mechanism exists purely to absorb benign duplicates. It is
    not a mechanism for deferring application cleanup or for re-dispatching
    responses to a replacement handler.
\end{remark}

\section{Pattern subscriptions}\label{sec:session_patterns}

Pattern subscriptions are indistinguishable from verbatim subscriptions on
the wire. The distinction exists only in how the local node decides which
topics to join.

\subsection{Discovery}

A node with a local pattern subscription discovers matching topics in two
complementary ways:

\begin{description}
    \item[Passive discovery] A received gossip carries the topic name (see
    section~\ref{sec:session_header_gossip}). If the received name matches a
    local pattern, the node creates an \emph{implicit topic} for it and
    attaches a subscriber. Passive discovery is automatic and requires no
    additional traffic.

    \item[Active discovery] A node that has just installed a new pattern
    subscription may broadcast a scout message
    (section~\ref{sec:session_scout}) with the pattern as the payload. Other
    nodes respond with gossip messages for any of their topics that match,
    which drives passive discovery.
\end{description}

\subsection{Substitutions}

When a pattern matches a name, the pattern's wildcards $(\texttt{*},\ \texttt{>})$
are paired with the segments they matched. The result is an ordered list of
\emph{substitutions}, each recording the matched substring together with its
ordinal position in the pattern.

\begin{remark}
    For example, the pattern \texttt{ins/*/data/>} matching the name
    \texttt{ins/0/data/foo/456} yields three substitutions:
    \verb|"0"|, \verb|"foo"|, and \verb|"456"|, at ordinals $0$, $1$, and $1$
    respectively (the second and third substitutions share an ordinal
    because they are both consumed by the trailing \texttt{>} token).
\end{remark}

The substitutions are computed once, when the pattern first matches the
topic, and are made available to the receiving application alongside each
arriving message. No string re-matching is performed per message.

\subsection{Implicit topic lifecycle}

An implicit topic behaves like any other topic for arbitration and message
delivery, with two differences:

\begin{itemize}
    \item It does not emit gossip of its own. Its CRDT state is updated by
    incoming gossip, by inline gossip from incoming application messages, or
    by scout-triggered passive discovery.
    \item It is automatically retired after the retention window elapses
    without any activity (no messages, no gossips).
\end{itemize}

The retention window should be much longer than the worst-case periodic
gossip interval so that a silent but alive topic is not prematurely
retired. A recommended default is ten minutes; see
section~\ref{sec:session_parameters}.

\subsection{Publisher invariant}

A publisher's topic name shall not contain any pattern tokens. A resolved
publisher name that contains \texttt{*} or \texttt{>} is invalid and shall
be rejected by the session layer at publisher creation time. This invariant
ensures that every topic on the wire has a deterministic, unambiguous name.

\section{Parameters and sizing}\label{sec:session_parameters}

The following parameters govern the wire behaviour of the session layer. Values
marked \emph{normative} are required for interoperability; values marked
\emph{recommended} may be adjusted by compliant implementations as long as the
CRDT invariants in section~\ref{sec:session_crdt} and the transport contract in
section~\ref{sec:transport} remain satisfied. Nodes that use different
recommended values remain interoperable with one another.

{\footnotesize
\begin{CyphalSimpleTable}[wide]{Session-layer parameters}{|l l l X|}
    Parameter                          & Unit  & Value              & Notes \\
    Session header size                & bytes & 24                 & Normative; section~\ref{sec:session_headers}. \\
    Maximum topic name length          & bytes & 200                & Normative; UTF-8, after normalization. \\
    Subject-ID modulus, 16-bit width   & --    & 57203              & Normative; section~\ref{sec:session_modulus}. \\
    Subject-ID modulus, 23-bit width   & --    & 8378431            & Normative; section~\ref{sec:session_modulus}. \\
    Subject-ID modulus, 32-bit width   & --    & 4294954663         & Normative; reserved for future 32-bit transports. \\
    Eviction counter width             & bits  & 32                 & Normative; pinned range starts at $\text{UINT32\_MAX} - 8191 = \text{0xFFFFE000}$. \\
    Message tag width                  & bits  & 64                 & Normative; random-initialized, permitted to wrap. \\
    Response sequence number width     & bits  & 48                 & Normative; monotonic per request, starts at zero. \\
    Log-age range                      & --    & $[-1,\ 35]$        & Normative; $-1$ encodes an age below one second. \\
    Default gossip period              & s     & $\sim 5$           & Recommended; randomized by $\pm 1/8$ of the period. \\
    Broadcast-to-shard ratio           & --    & $\sim 1 : 10$      & Recommended; one in $N$ periodic gossips goes on the broadcast subject. \\
    Urgent gossip delay window $W_u$   & ms    & $\sim 10$          & Recommended; see section~\ref{sec:session_gossip_delays}. \\
    Startup gossip delay window $W_s$  & ms    & tens to hundreds   & Recommended; see section~\ref{sec:session_gossip_delays}. \\
    Reliable ACK baseline timeout      & ms    & $\sim 16$          & Recommended; retransmissions use exponential backoff. \\
    Session lifetime (soft state)      & s     & $\sim 60$          & Recommended; association and deduplication entries expire after this. \\
    Implicit topic retention window    & s     & $\sim 600$         & Recommended; inactive implicit topics are retired after this interval. \\
    Association slack limit            & --    & $2$                & Recommended; consecutive missed acks before an association is forgotten. \\
\end{CyphalSimpleTable}
}

Auxiliary material that does not fit in the main body of the chapter is
collected in the appendices. Appendix~\ref{sec:appendix_eviction_solver}
gives a constant-time diagnostic solver for recovering the eviction counter
from a subject-ID. Appendix~\ref{sec:appendix_core_tla} reproduces the
complete TLA\textsuperscript{+} model of the allocation CRDT that this
chapter paraphrases; appendix~\ref{sec:appendix_proof} contains the
convergence proof. Appendix~\ref{sec:appendix_rapidhash} contains the pinned
rapidhash reference source.


\section{Interoperation with Cyphal/CAN v1.0}\label{sec:session_v10_interop}

Cyphal v1.1 nodes interoperate with Cyphal/CAN v1.0 nodes on the same bus.
The mechanism is transport-level rather than session-level: a v1.1 node sending a
best-effort message on a pinned topic in the range $[0, 8191]$ strips its
session-layer header and emits the payload as a legacy 13-bit frame, which v1.0 nodes
parse without modification.
See section~\ref{sec:transport_can_format_selection} and
section~\ref{sec:transport_can_v10_header_reconstruction} for the normative details of
format selection and header reconstruction.

v1.1 gossip, scout, RPC, and reliable traffic are confined to the 16-bit CAN ID format,
which v1.0 nodes silently discard at their bit~7 check.
v1.0 nodes therefore see only the subset of v1.1 traffic that lies on pinned subject-IDs
and is sent best-effort; everything else is invisible to them and does not affect them.

Pattern subscriptions cannot discover v1.0-only topics because v1.0 nodes emit neither
gossip nor scout messages. Applications that need to interact with v1.0 endpoints
shall use verbatim pinned subscriptions on the corresponding subject-IDs.
The session layer does not bridge v1.0 RPC: v1.0 service transfers use the
service-ID mechanism of the transport layer, which the session layer does not observe.

\section{Security and threat model}\label{sec:session_security}

The session layer assumes that all nodes on a Cyphal network are trusted.
The protocol is designed to remain correct under network faults, packet loss,
reordering, partial partitioning, and resource starvation, but it is not designed to
remain correct under adversarial behavior by a participating node.
A malicious node that is permitted to emit messages on the network can influence
topic allocation, inject gossip, or send acknowledgements on behalf of others.

Security features -- authentication, encryption, authorization, signed pinning -- may
be introduced as optional extensions once the core protocol has stabilized.
They are not specified here.
